% \documentclass[11pt,respuestas,a4]{aleph-examen}
\documentclass[11pt,a5]{aleph-examen}

% -- Paquetes extra
\usepackage{aleph-comandos}
\usepackage{systeme}
\input{formato_ipynb}

% -- Datos
\institucion{Facultad de Ciencias Exactas, Naturales y Ambientales}
\carrera{Catalogo STEM}
\asignatura{Algebra lineal}
\tema{Recuperación 1}
\autor{Andres Merino}
\fecha{Periodo 2026-1}

\logouno[0.16\textwidth]{Logos/logoPUCE_04_ac}
\definecolor{colortext}{HTML}{1957d1}
\definecolor{colordef}{HTML}{1957d1}
\fuente{montserrat}


\begin{document}

\encabezado

\section*{Ejercicios Matrices}

\begin{preguntas}

%%%%%%%%%%%%%%%%%%%%%%%%%%%%%%%%%%%%%%%%%%
%%%%%%%%%%%%%%%%%%%%%%%%%%%%%%%%%%%%%%%%%%
%%%%%%%%%%%%%%%%%%%%%%%%%%%%%%%%%%%%%%%%%%
%%%%%%%%%%%%%%%%%%%%%%%%%%%%%%%%%%%%%%%%%%
\item
Sea \( \alpha \in \mathbb{R} \) y considere la matriz
\[
    A=\begin{pmatrix}
        1 & 2 & 0 \\
        0 & 1 & 1 \\
        2 & \alpha & 3
    \end{pmatrix}.
\]

\begin{enumerate}
    \item Utilizando expansión por cofactores en la primera columna, calcule \( \det(A) \).
    \item Determine los valores de \( \alpha \) para los cuales la matriz es no singular.
    \item Calcule \( A^{-1} \) mediante reducción por filas, para uno de los valores de $\alpha$ obtenidos en el ejercicio anterior.
\end{enumerate}

%%%%%%%%%%%%%%%%%%%%%%%%%%%%%%%%%%%%%%%%%%
\item
Sean \( \alpha \in \mathbb{R} \) y las matrices
\[
    A=\begin{pmatrix}
        1 & \alpha \\
        0 & 1
    \end{pmatrix},
    \qquad
    B=\begin{pmatrix}
        1 & 1 \\
        0 & 1
    \end{pmatrix}.
\]

\begin{enumerate}
    \item Calcule \( AB \) y \( BA \).
    \item Determine los valores de \( \alpha \) para los cuales \( A \) y \( B \) conmutan, es decir, \( AB=BA \).
\end{enumerate}

\end{preguntas}

\section*{Ejercicios Sistemas de Ecuaciones}

\begin{preguntas}
%%%%%%%%%%%%%%%%%%%%%%%%%%%%%%%%%%%%%%%%%%
\item
Sea \( \alpha \in \mathbb{R} \) y considere el sistema lineal
\[
\systeme[xyz]{
x+y+z=2{,},
2x+3y+\alpha z=5{,},
x+2y+{(\alpha-1)} z=3{.}
}
\]

\begin{enumerate}
    \item Escriba la matriz ampliada del sistema.
    \item Aplique reducción de Gauss-Jordan.
    \item Determine el valor de \( \alpha \) para el cual el sistema tiene infinitas soluciones.
    \item Para dicho valor, escriba el conjunto solución en forma paramétrica.
\end{enumerate}

%%%%%%%%%%%%%%%%%%%%%%%%%%%%%%%%%%%%%%%%%%
\item
Considere el sistema lineal de ecuaciones
\[
\systeme{
x+y+z=6,
x+2y+3z=14,
2x+y+z=9
}
\]

\begin{enumerate}
    \item Escriba la matriz de coeficientes y calcule su determinante.
    \item Verifique que la matriz es invertible.
    \item Sabiendo que
    \[
        A^{-1}=
        \begin{pmatrix}
            1 & 0 & -1 \\
            1 & -1 & 2 \\
            -1 & 1 & 0
        \end{pmatrix},
    \]
    utilice la matriz inversa para resolver el sistema.
\end{enumerate}


\end{preguntas}

\end{document}
